\documentclass[11pt,a4paper]{article}

\usepackage{amsmath,amssymb,amsfonts}
\usepackage{bm}
\usepackage{geometry}
\usepackage{booktabs}
\usepackage{array}
\usepackage{hyperref}
\usepackage{cleveref}
\usepackage{siunitx}
\usepackage{xcolor}
\usepackage{fancyhdr}

\geometry{margin=2.5cm, headheight=14pt}

\hypersetup{
    colorlinks=true,
    linkcolor=blue!70!black,
    citecolor=green!50!black,
    urlcolor=blue!70!black
}

\pagestyle{fancy}
\fancyhf{}
\rhead{INS Error-State EKF}
\lhead{Technical Reference}
\rfoot{Page \thepage}

% Math commands
\newcommand{\mat}[1]{\mathbf{#1}}
\newcommand{\vect}[1]{\boldsymbol{#1}}
\newcommand{\quat}[1]{\mathbf{#1}}
\newcommand{\skewmat}[1]{\left[#1 \times\right]}
\newcommand{\Rn}{\mathbb{R}}

\title{\textbf{Inertial Navigation System}\\[0.5em]
\Large Error-State Extended Kalman Filter\\[0.3em]
\large Based on Farrell's State Augmentation with NBK Error Model}
\author{Technical Reference Document}
\date{\today}

\begin{document}

\maketitle
\tableofcontents
\newpage

%==============================================================================
\section{Introduction}
%==============================================================================

This document derives an error-state Extended Kalman Filter (ES-EKF) for strapdown inertial navigation. The formulation follows J.A. Farrell's state augmentation methodology with the NBK stochastic error model, where parameters are determined via Allan variance analysis.

\subsection{Error-State EKF Methodology}

The ES-EKF maintains two parallel representations:

\begin{enumerate}
    \item \textbf{Nominal state} $\mat{x}$: The best estimate of the true state, propagated by nonlinear mechanization equations. Uses quaternion for attitude (4 parameters).
    
    \item \textbf{Error state} $\delta\mat{x}$: Small perturbations from the nominal state, estimated by the Kalman filter. Uses a minimal 3-parameter attitude error representation.
\end{enumerate}

This separation is essential because:
\begin{itemize}
    \item Quaternions provide singularity-free attitude representation but have 4 parameters for 3 degrees of freedom
    \item The unit-norm constraint $\|\quat{q}\| = 1$ makes direct quaternion estimation problematic
    \item A 3-parameter attitude error respects the 3 DOF of rotation and yields well-conditioned covariance matrices
\end{itemize}

\subsection{Notation}

\begin{center}
\begin{tabular}{ll}
    \toprule
    Symbol & Description \\
    \midrule
    $(\cdot)^b$ & Body frame quantity \\
    $(\cdot)^n$ & Navigation frame (NED) quantity \\
    $\mat{C}_b^n$ & Direction cosine matrix, body to navigation \\
    $\quat{q}_b^n$ & Unit quaternion, body to navigation (Hamilton, scalar-last) \\
    $\delta(\cdot)$ & Error in quantity \\
    $\tilde{(\cdot)}$ & Measured (noisy) quantity \\
    $\skewmat{\vect{a}}$ & Skew-symmetric matrix of vector $\vect{a}$ \\
    $\otimes$ & Quaternion multiplication \\
    \bottomrule
\end{tabular}
\end{center}

\subsection{Quaternion Convention}

Throughout this document, quaternions use the \textbf{Hamilton convention with scalar-last ordering}:
\begin{equation}
    \quat{q} = \begin{bmatrix} q_x \\ q_y \\ q_z \\ q_w \end{bmatrix} = \begin{bmatrix} \vect{q}_v \\ q_w \end{bmatrix}
\end{equation}
where $\vect{q}_v \in \Rn^3$ is the vector part and $q_w \in \Rn$ is the scalar part. The quaternion satisfies $\|\quat{q}\|^2 = q_x^2 + q_y^2 + q_z^2 + q_w^2 = 1$.

%==============================================================================
\section{State Definitions}
%==============================================================================

\subsection{Nominal State (16 parameters)}

The nominal state vector propagated by the mechanization:
\begin{equation}
    \mat{x} = \begin{bmatrix} \mat{r} \\ \mat{v}^n \\ \quat{q}_b^n \\ \mat{b}_a \\ \mat{b}_g \end{bmatrix} \in \Rn^{16}
\end{equation}

\begin{center}
\begin{tabular}{clcc}
    \toprule
    Index & Component & Symbol & Dimension \\
    \midrule
    0--2 & Geodetic position & $\mat{r} = [\phi, \lambda, h]^T$ & 3 \\
    3--5 & NED velocity & $\mat{v}^n = [v_N, v_E, v_D]^T$ & 3 \\
    6--9 & Attitude quaternion & $\quat{q}_b^n = [q_x, q_y, q_z, q_w]^T$ & 4 \\
    10--12 & Accelerometer bias & $\mat{b}_a = [b_{ax}, b_{ay}, b_{az}]^T$ & 3 \\
    13--15 & Gyroscope bias & $\mat{b}_g = [b_{gx}, b_{gy}, b_{gz}]^T$ & 3 \\
    \bottomrule
\end{tabular}
\end{center}

\subsection{Error State (21 parameters)}

The error state estimated by the Kalman filter:
\begin{equation}
    \delta\mat{x} = \begin{bmatrix} \delta\mat{r} \\ \delta\mat{v}^n \\ \vect{\theta} \\ \delta\mat{b}_a \\ \delta\mat{b}_g \end{bmatrix} \in \Rn^{21}
\end{equation}

\begin{center}
\begin{tabular}{clcc}
    \toprule
    Index & Component & Symbol & Dimension \\
    \midrule
    0--2 & Position error & $\delta\mat{r} = [\delta\phi, \delta\lambda, \delta h]^T$ & 3 \\
    3--5 & Velocity error & $\delta\mat{v}^n = [\delta v_N, \delta v_E, \delta v_D]^T$ & 3 \\
    6--8 & Attitude error & $\vect{\theta} = [\theta_N, \theta_E, \theta_D]^T$ & 3 \\
    9--11 & Accel bias error & $\delta\mat{b}_a$ & 3 \\
    12--14 & Gyro bias error & $\delta\mat{b}_g$ & 3 \\
    15--17 & Accel bias (stochastic) & $\mat{b}_a^z$ & 3 \\
    18--20 & Gyro bias (stochastic) & $\mat{b}_g^z$ & 3 \\
    \bottomrule
\end{tabular}
\end{center}

\textbf{Note}: Following Farrell's state augmentation, the bias states appear twice: as errors in the nominal bias estimates ($\delta\mat{b}_a$, $\delta\mat{b}_g$) and as stochastic states driven by rate random walk ($\mat{b}_a^z$, $\mat{b}_g^z$). For the simplified NBK model where biases are modeled as pure random walks, these can be combined, yielding an error state in $\Rn^{15}$. This document presents the full augmented form.

\textbf{Simplified 15-state formulation}: If we model biases directly as random walk without separate deterministic error:
\begin{equation}
    \delta\mat{x} = \begin{bmatrix} \delta\mat{r} \\ \delta\mat{v}^n \\ \vect{\theta} \\ \mat{b}_a \\ \mat{b}_g \end{bmatrix} \in \Rn^{15}
\end{equation}
This is the formulation used in the remainder of this document.

\subsection{Attitude Error and Quaternion Correction}

The attitude error vector $\vect{\theta} \in \Rn^3$ represents a small rotation. The error quaternion is constructed as:
\begin{equation}
    \delta\quat{q} = \begin{bmatrix} \frac{1}{2}\vect{\theta} \\ 1 \end{bmatrix} \cdot \frac{1}{\sqrt{1 + \frac{1}{4}\|\vect{\theta}\|^2}}
    \label{eq:error_quat}
\end{equation}

For small errors, this simplifies to:
\begin{equation}
    \delta\quat{q} \approx \begin{bmatrix} \frac{1}{2}\vect{\theta} \\ 1 \end{bmatrix}
\end{equation}

The corrected quaternion after a Kalman update:
\begin{equation}
    \quat{q}^+ = \delta\quat{q} \otimes \quat{q}^-
    \label{eq:quat_correction}
\end{equation}

After correction, the quaternion must be renormalized: $\quat{q}^+ \leftarrow \quat{q}^+ / \|\quat{q}^+\|$.

%==============================================================================
\section{Quaternion Operations}
%==============================================================================

\subsection{Quaternion Product}

For quaternions $\quat{p} = [\vect{p}_v^T, p_w]^T$ and $\quat{q} = [\vect{q}_v^T, q_w]^T$:
\begin{equation}
    \quat{p} \otimes \quat{q} = \begin{bmatrix}
        p_w \vect{q}_v + q_w \vect{p}_v + \vect{p}_v \times \vect{q}_v \\
        p_w q_w - \vect{p}_v \cdot \vect{q}_v
    \end{bmatrix}
\end{equation}

In matrix form:
\begin{equation}
    \quat{p} \otimes \quat{q} = \mat{L}(\quat{p}) \quat{q} = \mat{R}(\quat{q}) \quat{p}
\end{equation}

where:
\begin{align}
    \mat{L}(\quat{p}) &= \begin{bmatrix}
        p_w \mat{I}_3 + \skewmat{\vect{p}_v} & \vect{p}_v \\
        -\vect{p}_v^T & p_w
    \end{bmatrix} \\
    \mat{R}(\quat{q}) &= \begin{bmatrix}
        q_w \mat{I}_3 - \skewmat{\vect{q}_v} & \vect{q}_v \\
        -\vect{q}_v^T & q_w
    \end{bmatrix}
\end{align}

\subsection{Quaternion to DCM}

The direction cosine matrix from quaternion:
\begin{equation}
    \mat{C}_b^n = (q_w^2 - \vect{q}_v^T\vect{q}_v)\mat{I}_3 + 2\vect{q}_v\vect{q}_v^T + 2q_w\skewmat{\vect{q}_v}
\end{equation}

Expanded:
\begin{equation}
    \mat{C}_b^n = \begin{bmatrix}
        1-2(q_y^2+q_z^2) & 2(q_xq_y-q_zq_w) & 2(q_xq_z+q_yq_w) \\
        2(q_xq_y+q_zq_w) & 1-2(q_x^2+q_z^2) & 2(q_yq_z-q_xq_w) \\
        2(q_xq_z-q_yq_w) & 2(q_yq_z+q_xq_w) & 1-2(q_x^2+q_y^2)
    \end{bmatrix}
    \label{eq:quat_to_dcm}
\end{equation}

\subsection{Quaternion Kinematics}

The quaternion rate equation:
\begin{equation}
    \dot{\quat{q}}_b^n = \frac{1}{2} \quat{q}_b^n \otimes \quat{\omega}_{nb}^b
    \label{eq:quat_kinematics}
\end{equation}

where $\quat{\omega}_{nb}^b = [\vect{\omega}_{nb}^{b\,T}, 0]^T$ is the angular velocity as a pure quaternion, and:
\begin{equation}
    \vect{\omega}_{nb}^b = \vect{\omega}_{ib}^b - \mat{C}_n^b \vect{\omega}_{in}^n
\end{equation}

In matrix form:
\begin{equation}
    \dot{\quat{q}}_b^n = \frac{1}{2} \mat{\Omega}(\vect{\omega}_{nb}^b) \quat{q}_b^n
\end{equation}

with:
\begin{equation}
    \mat{\Omega}(\vect{\omega}) = \begin{bmatrix}
        -\skewmat{\vect{\omega}} & \vect{\omega} \\
        -\vect{\omega}^T & 0
    \end{bmatrix}
\end{equation}

\subsection{Skew-Symmetric Matrix}

For $\vect{a} = [a_1, a_2, a_3]^T$:
\begin{equation}
    \skewmat{\vect{a}} = \begin{bmatrix}
        0 & -a_3 & a_2 \\
        a_3 & 0 & -a_1 \\
        -a_2 & a_1 & 0
    \end{bmatrix}
\end{equation}

Property: $\skewmat{\vect{a}}\vect{b} = \vect{a} \times \vect{b}$

%==============================================================================
\section{Continuous-Time Error Dynamics}
%==============================================================================

The linearized error-state dynamics:
\begin{equation}
    \delta\dot{\mat{x}} = \mat{F}\,\delta\mat{x} + \mat{G}\mat{w}
    \label{eq:error_dynamics}
\end{equation}

where $\mat{w}$ is white noise with spectral density $\mat{Q}_c$.

\subsection{System Matrix Structure}

The $15 \times 15$ system matrix:
\begin{equation}
    \mat{F} = \begin{bmatrix}
        \mat{F}_{rr} & \mat{F}_{rv} & \mat{0}_{3\times3} & \mat{0}_{3\times3} & \mat{0}_{3\times3} \\
        \mat{F}_{vr} & \mat{F}_{vv} & \mat{F}_{v\theta} & \mat{C}_b^n & \mat{0}_{3\times3} \\
        \mat{F}_{\theta r} & \mat{F}_{\theta v} & \mat{F}_{\theta\theta} & \mat{0}_{3\times3} & -\mat{C}_b^n \\
        \mat{0}_{3\times3} & \mat{0}_{3\times3} & \mat{0}_{3\times3} & \mat{0}_{3\times3} & \mat{0}_{3\times3} \\
        \mat{0}_{3\times3} & \mat{0}_{3\times3} & \mat{0}_{3\times3} & \mat{0}_{3\times3} & \mat{0}_{3\times3}
    \end{bmatrix}
    \label{eq:F_matrix}
\end{equation}

State indices: $\delta\mat{r}$ (0:3), $\delta\mat{v}$ (3:6), $\vect{\theta}$ (6:9), $\mat{b}_a$ (9:12), $\mat{b}_g$ (12:15).

\subsection{Position Error Dynamics}

Position rate depends on velocity:
\begin{equation}
    \mat{F}_{rr} = \begin{bmatrix}
        0 & 0 & \dfrac{-v_N}{(R_M+h)^2} \\[1em]
        \dfrac{v_E \tan\phi \sec\phi}{R_N+h} & 0 & \dfrac{-v_E \sec\phi}{(R_N+h)^2} \\[1em]
        0 & 0 & 0
    \end{bmatrix}
\end{equation}

\begin{equation}
    \mat{F}_{rv} = \begin{bmatrix}
        \dfrac{1}{R_M+h} & 0 & 0 \\[1em]
        0 & \dfrac{\sec\phi}{R_N+h} & 0 \\[1em]
        0 & 0 & -1
    \end{bmatrix}
\end{equation}

\subsection{Velocity Error Dynamics}

The attitude error couples to velocity via specific force:
\begin{equation}
    \mat{F}_{v\theta} = -\skewmat{\mat{f}^n} = \begin{bmatrix}
        0 & f_D & -f_E \\
        -f_D & 0 & f_N \\
        f_E & -f_N & 0
    \end{bmatrix}
\end{equation}

where $\mat{f}^n = \mat{C}_b^n (\tilde{\mat{f}}^b - \mat{b}_a)$ is specific force in navigation frame.

The accelerometer bias couples directly:
\begin{equation}
    \mat{F}_{v,b_a} = \mat{C}_b^n
\end{equation}

\subsection{Attitude Error Dynamics}

\begin{equation}
    \mat{F}_{\theta\theta} = -\skewmat{\vect{\omega}_{in}^n}
\end{equation}

The gyroscope bias couples as:
\begin{equation}
    \mat{F}_{\theta,b_g} = -\mat{C}_b^n
\end{equation}

\subsection{Navigation Frame Rotation Rates}

Earth rate in navigation frame:
\begin{equation}
    \vect{\omega}_{ie}^n = \begin{bmatrix}
        \omega_e \cos\phi \\ 0 \\ -\omega_e \sin\phi
    \end{bmatrix}, \quad \omega_e = \SI{7.292115e-5}{\radian\per\second}
\end{equation}

Transport rate:
\begin{equation}
    \vect{\omega}_{en}^n = \begin{bmatrix}
        \dfrac{v_E}{R_N + h} \\[1em]
        \dfrac{-v_N}{R_M + h} \\[1em]
        \dfrac{-v_E \tan\phi}{R_N + h}
    \end{bmatrix}
\end{equation}

Total rotation:
\begin{equation}
    \vect{\omega}_{in}^n = \vect{\omega}_{ie}^n + \vect{\omega}_{en}^n
\end{equation}

\subsection{Earth Ellipsoid Parameters}

\begin{align}
    R_0 &= \SI{6378137.0}{\meter} & &\text{(semi-major axis)} \\
    e^2 &= 0.00669437999014 & &\text{(first eccentricity squared)} \\
    R_M &= \frac{R_0(1-e^2)}{(1 - e^2\sin^2\phi)^{3/2}} & &\text{(meridian radius)} \\
    R_N &= \frac{R_0}{\sqrt{1 - e^2\sin^2\phi}} & &\text{(prime vertical radius)}
\end{align}

%==============================================================================
\section{IMU Stochastic Error Model (NBK)}
%==============================================================================

\subsection{Allan Variance Parameters}

From Farrell's Table 1, the dominant error sources characterized by Allan variance:

\begin{center}
\begin{tabular}{llll}
    \toprule
    \textbf{Parameter} & \textbf{Symbol} & \textbf{Accel Units} & \textbf{Gyro Units} \\
    \midrule
    Velocity/Angle Random Walk & $N$ & \si{\meter\per\second\per\sqrt{\hour}} & \si{\degree\per\sqrt{\hour}} \\
    Bias Instability & $B$ & \si{\micro g} & \si{\degree\per\hour} \\
    Accel/Rate Random Walk & $K$ & \si{\micro g\per\sqrt{\second}} & \si{\degree\per\hour\per\sqrt{\second}} \\
    \bottomrule
\end{tabular}
\end{center}

\subsection{Allan Variance Expressions}

\begin{align}
    \text{Random Walk:} \quad \sigma^2(\tau) &= \frac{N^2}{\tau} \\
    \text{Bias Instability:} \quad \sigma^2(\tau) &= \frac{2B^2 \ln 2}{\pi} \quad \text{(at minimum)} \\
    \text{Rate Random Walk:} \quad \sigma^2(\tau) &= \frac{K^2 \tau}{3}
\end{align}

\subsection{Stochastic Model}

For the NBK model, sensor biases follow a random walk driven by rate random walk $K$:
\begin{align}
    \dot{\mat{b}}_a &= \mat{w}_{K_a}, \quad E[\mat{w}_{K_a}\mat{w}_{K_a}^T] = \text{diag}(K_{ax}^2, K_{ay}^2, K_{az}^2)\,\delta(t-\tau) \\
    \dot{\mat{b}}_g &= \mat{w}_{K_g}, \quad E[\mat{w}_{K_g}\mat{w}_{K_g}^T] = \text{diag}(K_{gx}^2, K_{gy}^2, K_{gz}^2)\,\delta(t-\tau)
\end{align}

The measurement model includes additive white noise:
\begin{align}
    \tilde{\mat{f}}^b &= \mat{f}^b + \mat{b}_a + \vect{\eta}_a \\
    \tilde{\vect{\omega}}^b &= \vect{\omega}^b + \mat{b}_g + \vect{\eta}_g
\end{align}

where $\vect{\eta}_a$, $\vect{\eta}_g$ are white noise with PSDs $N_a^2$, $N_g^2$.

%==============================================================================
\section{Noise Input Matrix and Spectral Density}
%==============================================================================

\subsection{Noise Input Matrix}

The $15 \times 12$ matrix mapping noise to error states:
\begin{equation}
    \mat{G} = \begin{bmatrix}
        \mat{0}_{3\times3} & \mat{0}_{3\times3} & \mat{0}_{3\times3} & \mat{0}_{3\times3} \\
        \mat{C}_b^n & \mat{0}_{3\times3} & \mat{0}_{3\times3} & \mat{0}_{3\times3} \\
        \mat{0}_{3\times3} & -\mat{C}_b^n & \mat{0}_{3\times3} & \mat{0}_{3\times3} \\
        \mat{0}_{3\times3} & \mat{0}_{3\times3} & \mat{I}_{3} & \mat{0}_{3\times3} \\
        \mat{0}_{3\times3} & \mat{0}_{3\times3} & \mat{0}_{3\times3} & \mat{I}_{3}
    \end{bmatrix}
    \label{eq:G_matrix}
\end{equation}

Column assignments:
\begin{center}
\begin{tabular}{cl}
    \toprule
    Columns & Noise Source \\
    \midrule
    0--2 & Accelerometer white noise $\vect{\eta}_a$ \\
    3--5 & Gyroscope white noise $\vect{\eta}_g$ \\
    6--8 & Accelerometer rate random walk $\mat{w}_{K_a}$ \\
    9--11 & Gyroscope rate random walk $\mat{w}_{K_g}$ \\
    \bottomrule
\end{tabular}
\end{center}

\subsection{Continuous Spectral Density}

The $12 \times 12$ spectral density:
\begin{equation}
    \mat{Q}_c = \text{diag}\left(
        N_{ax}^2, N_{ay}^2, N_{az}^2,
        N_{gx}^2, N_{gy}^2, N_{gz}^2,
        K_{ax}^2, K_{ay}^2, K_{az}^2,
        K_{gx}^2, K_{gy}^2, K_{gz}^2
    \right)
    \label{eq:Qc}
\end{equation}

\subsection{Unit Conversions}

For accelerometer data in $g$ and gyroscope data in \si{\radian\per\second}:
\begin{align}
    N_a^{\text{SI}} &= N_a^{(g)} \cdot g_0, \quad g_0 = \SI{9.80665}{\meter\per\second\squared} \\
    K_a^{\text{SI}} &= K_a^{(g)} \cdot g_0
\end{align}

Gyroscope parameters in \si{\radian\per\second} require no conversion.

%==============================================================================
\section{Van Loan Discretization}
%==============================================================================

The Van Loan method computes the exact discrete-time state transition matrix $\vect{\Phi}$ and process noise covariance $\mat{Q}_d$ from the continuous system.

\subsection{Construction}

Form the $2n \times 2n$ matrix ($n = 15$):
\begin{equation}
    \mat{A} = \begin{bmatrix}
        -\mat{F} & \mat{G}\mat{Q}_c\mat{G}^T \\
        \mat{0}_{n\times n} & \mat{F}^T
    \end{bmatrix} \Delta t
    \label{eq:van_loan_A}
\end{equation}

\subsection{Matrix Exponential}

Compute:
\begin{equation}
    \mat{B} = \exp(\mat{A}) = \begin{bmatrix}
        \mat{B}_{11} & \mat{B}_{12} \\
        \mat{0} & \mat{B}_{22}
    \end{bmatrix}
\end{equation}

\subsection{Extract Discrete Quantities}

\begin{align}
    \vect{\Phi} &= \mat{B}_{22}^T \label{eq:Phi_extract} \\
    \mat{Q}_d &= \vect{\Phi} \mat{B}_{12} \label{eq:Qd_extract}
\end{align}

Symmetrize: $\mat{Q}_d \leftarrow \frac{1}{2}(\mat{Q}_d + \mat{Q}_d^T)$

\subsection{Process Noise Projection}

The $15 \times 15$ matrix $\mat{G}\mat{Q}_c\mat{G}^T$:
\begin{equation}
    \mat{G}\mat{Q}_c\mat{G}^T = \begin{bmatrix}
        \mat{0}_{3\times3} & & & & \\
        & \mat{C}_b^n\mat{Q}_{N_a}\mat{C}_n^b & & & \\
        & & \mat{C}_b^n\mat{Q}_{N_g}\mat{C}_n^b & & \\
        & & & \mat{Q}_{K_a} & \\
        & & & & \mat{Q}_{K_g}
    \end{bmatrix}
\end{equation}

where off-diagonal blocks are zero and:
\begin{align}
    \mat{Q}_{N_a} &= \text{diag}(N_{ax}^2, N_{ay}^2, N_{az}^2) \\
    \mat{Q}_{N_g} &= \text{diag}(N_{gx}^2, N_{gy}^2, N_{gz}^2) \\
    \mat{Q}_{K_a} &= \text{diag}(K_{ax}^2, K_{ay}^2, K_{az}^2) \\
    \mat{Q}_{K_g} &= \text{diag}(K_{gx}^2, K_{gy}^2, K_{gz}^2)
\end{align}

%==============================================================================
\section{Measurement Models}
%==============================================================================

All measurements relate to the error state via:
\begin{equation}
    \mat{z} = \mat{H}\,\delta\mat{x} + \vect{\nu}
\end{equation}

where $\vect{\nu}$ is measurement noise with covariance $\mat{R}$.

\subsection{GNSS Position}

Innovation:
\begin{equation}
    \mat{z}_{\text{pos}} = \mat{r}_{\text{GNSS}} - \hat{\mat{r}}_{\text{INS}}
\end{equation}

Observation matrix ($3 \times 15$):
\begin{equation}
    \mat{H}_{\text{pos}} = \begin{bmatrix}
        \mat{I}_{3} & \mat{0}_{3\times12}
    \end{bmatrix}
\end{equation}

\subsection{GNSS Velocity}

Innovation:
\begin{equation}
    \mat{z}_{\text{vel}} = \mat{v}_{\text{GNSS}}^n - \hat{\mat{v}}_{\text{INS}}^n
\end{equation}

Observation matrix ($3 \times 15$):
\begin{equation}
    \mat{H}_{\text{vel}} = \begin{bmatrix}
        \mat{0}_{3\times3} & \mat{I}_{3} & \mat{0}_{3\times9}
    \end{bmatrix}
\end{equation}

\subsection{Barometric Altitude}

From pressure $p$ in \si{\hecto\pascal}:
\begin{equation}
    h_{\text{baro}} = \frac{T_0}{L}\left[1 - \left(\frac{p}{p_0}\right)^{\frac{LR}{gM}}\right]
\end{equation}

with $T_0 = \SI{288.15}{\kelvin}$, $L = \SI{0.0065}{\kelvin\per\meter}$, $p_0 = \SI{1013.25}{\hecto\pascal}$.

Innovation:
\begin{equation}
    z_{\text{baro}} = h_{\text{baro}} - \hat{h}_{\text{INS}}
\end{equation}

Observation matrix ($1 \times 15$):
\begin{equation}
    \mat{H}_{\text{baro}} = \begin{bmatrix}
        0 & 0 & 1 & \mat{0}_{1\times12}
    \end{bmatrix}
\end{equation}

\subsection{Zero Velocity Update (ZUPT)}

When stationary (detected via gyro/accel magnitude):
\begin{equation}
    \mat{z}_{\text{ZUPT}} = \mat{0}_{3\times1} - \hat{\mat{v}}_{\text{INS}}^n
\end{equation}

Uses $\mat{H}_{\text{vel}}$.

\subsection{Magnetometer Heading}

The magnetometer provides heading observability by measuring Earth's magnetic field. A heading-only formulation avoids interference with accelerometer-derived tilt estimates.

\subsubsection{Magnetic Field Model}

The magnetometer measures the local magnetic field in the body frame:
\begin{equation}
    \tilde{\mat{m}}^b = \mat{C}_n^b \mat{m}^n + \mat{b}_m + \vect{\nu}_m
\end{equation}

where $\mat{m}^n$ is the reference field from a geomagnetic model (WMM/IGRF), $\mat{b}_m$ is hard-iron bias, and $\vect{\nu}_m$ is measurement noise.

The reference field in NED:
\begin{equation}
    \mat{m}^n = \begin{bmatrix} m_N \\ m_E \\ m_D \end{bmatrix} = \|\mat{m}\| \begin{bmatrix} \cos\delta\cos\iota \\ \sin\delta\cos\iota \\ \sin\iota \end{bmatrix}
\end{equation}

where $\delta$ is magnetic declination and $\iota$ is magnetic inclination (both from the geomagnetic model at current position).

\subsubsection{Heading Extraction}

Transform the calibrated magnetometer reading to the navigation frame:
\begin{equation}
    \hat{\mat{m}}^n = \mat{C}_b^n (\tilde{\mat{m}}^b - \hat{\mat{b}}_m)
\end{equation}

Extract magnetic heading:
\begin{equation}
    \psi_{\text{mag}} = \text{atan2}(\hat{m}_E^n, \hat{m}_N^n)
\end{equation}

The true heading is obtained by applying declination:
\begin{equation}
    \psi_{\text{true}} = \psi_{\text{mag}} + \delta
\end{equation}

\subsubsection{Innovation}

Compare magnetometer-derived heading to INS heading:
\begin{equation}
    z_{\psi} = \psi_{\text{mag}} - \hat{\psi}_{\text{INS}}
\end{equation}

The INS heading is extracted from the quaternion via the DCM:
\begin{equation}
    \hat{\psi}_{\text{INS}} = \text{atan2}(C_{12}, C_{11})
\end{equation}

where $C_{ij}$ are elements of $\mat{C}_b^n$.

\textbf{Angle wrapping}: The innovation must be wrapped to $[-\pi, \pi]$:
\begin{equation}
    z_{\psi} \leftarrow \text{atan2}(\sin z_{\psi}, \cos z_{\psi})
\end{equation}

\subsubsection{Observation Matrix}

The heading depends on the attitude error. For the linearization:
\begin{equation}
    \delta\psi = \frac{\partial \psi}{\partial \vect{\theta}} \vect{\theta}
\end{equation}

For level flight (small roll $\phi$ and pitch $\theta$), the heading error is dominated by yaw error:
\begin{equation}
    \mat{H}_{\psi}^{\text{approx}} = \begin{bmatrix}
        \mat{0}_{1\times6} & 0 & 0 & 1 & \mat{0}_{1\times6}
    \end{bmatrix}
\end{equation}

For the general case with arbitrary attitude, the full Jacobian is required. Let $\mat{C} = \mat{C}_b^n$:
\begin{equation}
    \frac{\partial \psi}{\partial \vect{\theta}} = \frac{1}{C_{11}^2 + C_{12}^2} \begin{bmatrix}
        C_{11}\dfrac{\partial C_{12}}{\partial \theta_N} - C_{12}\dfrac{\partial C_{11}}{\partial \theta_N} \\[1em]
        C_{11}\dfrac{\partial C_{12}}{\partial \theta_E} - C_{12}\dfrac{\partial C_{11}}{\partial \theta_E} \\[1em]
        C_{11}\dfrac{\partial C_{12}}{\partial \theta_D} - C_{12}\dfrac{\partial C_{11}}{\partial \theta_D}
    \end{bmatrix}^T
\end{equation}

The DCM perturbation due to attitude error:
\begin{equation}
    \delta\mat{C}_b^n = -\skewmat{\vect{\theta}}\mat{C}_b^n
\end{equation}

yields:
\begin{align}
    \delta C_{11} &= \theta_D C_{21} - \theta_E C_{31} \\
    \delta C_{12} &= \theta_D C_{22} - \theta_E C_{32}
\end{align}

Therefore:
\begin{equation}
    \frac{\partial \psi}{\partial \vect{\theta}} = \frac{1}{C_{11}^2 + C_{12}^2} \begin{bmatrix}
        0 \\
        -(C_{11}C_{32} - C_{12}C_{31}) \\
        C_{11}C_{22} - C_{12}C_{21}
    \end{bmatrix}^T
\end{equation}

The full observation matrix ($1 \times 15$):
\begin{equation}
    \mat{H}_{\psi} = \begin{bmatrix}
        \mat{0}_{1\times6} & \dfrac{\partial \psi}{\partial \vect{\theta}} & \mat{0}_{1\times6}
    \end{bmatrix}
    \label{eq:H_mag}
\end{equation}

\subsubsection{Measurement Noise}

Typical magnetometer noise variance:
\begin{equation}
    R_{\psi} = \sigma_{\psi}^2
\end{equation}

where $\sigma_{\psi}$ depends on sensor quality and magnetic environment (typically $\SI{1}{\degree}$ to $\SI{5}{\degree}$).

\subsubsection{Magnetic Disturbance Detection}

Reject magnetometer updates when the field magnitude deviates significantly from the expected value:
\begin{equation}
    \text{if } \left| \|\tilde{\mat{m}}^b\| - \|\mat{m}^n\| \right| > \epsilon_m \quad \Rightarrow \quad \text{skip update}
\end{equation}

Typical threshold: $\epsilon_m \approx 0.1 \|\mat{m}^n\|$ (10\% deviation).

\subsubsection{Update Rate}

Magnetometer updates are typically applied at a lower rate than IMU propagation:
\begin{itemize}
    \item IMU: 100--400 Hz
    \item Magnetometer: 1--20 Hz
\end{itemize}

This reduces computational load and limits the influence of high-frequency magnetic noise.

%==============================================================================
\section{EKF Algorithm}
%==============================================================================

\subsection{Initialization}

\begin{enumerate}
    \item Set initial position from GNSS or known location
    \item Initialize velocity to zero (static alignment) or from GNSS
    \item Compute initial quaternion from accelerometer (gravity) and optionally magnetometer
    \item Set initial biases to zero or from calibration
    \item Initialize $\mat{P}_0$ with appropriate uncertainties
\end{enumerate}

\subsection{Prediction Step}

For each IMU sample at rate $f_s = 1/\Delta t$:

\begin{enumerate}
    \item \textbf{Correct measurements}:
    \begin{align}
        \mat{f}^b &= \tilde{\mat{f}}^b - \hat{\mat{b}}_a \\
        \vect{\omega}^b &= \tilde{\vect{\omega}}^b - \hat{\mat{b}}_g
    \end{align}
    
    \item \textbf{Propagate nominal state} via strapdown mechanization:
    \begin{align}
        \quat{q}_{k+1} &= \quat{q}_k \otimes \Delta\quat{q}(\vect{\omega}_{nb}^b, \Delta t) \\
        \mat{v}_{k+1}^n &= \mat{v}_k^n + (\mat{C}_b^n \mat{f}^b + \mat{g}^n - (2\vect{\omega}_{ie}^n + \vect{\omega}_{en}^n) \times \mat{v}_k^n)\Delta t \\
        \mat{r}_{k+1} &= \mat{r}_k + \mat{F}_{rv}^{-1}\mat{v}_k^n \Delta t
    \end{align}
    
    \item \textbf{Build $\mat{F}$, $\mat{G}$, $\mat{Q}_c$} at current state
    
    \item \textbf{Van Loan discretization} $\rightarrow$ $\vect{\Phi}$, $\mat{Q}_d$
    
    \item \textbf{Propagate covariance}:
    \begin{equation}
        \mat{P}_{k+1}^- = \vect{\Phi}\mat{P}_k^+\vect{\Phi}^T + \mat{Q}_d
    \end{equation}
\end{enumerate}

\subsection{Update Step}

When measurement available:

\begin{enumerate}
    \item \textbf{Compute innovation}:
    \begin{equation}
        \mat{y} = \mat{z} - \mat{H}\hat{\mat{x}}^-
    \end{equation}
    
    \item \textbf{Innovation covariance}:
    \begin{equation}
        \mat{S} = \mat{H}\mat{P}^-\mat{H}^T + \mat{R}
    \end{equation}
    
    \item \textbf{Kalman gain}:
    \begin{equation}
        \mat{K} = \mat{P}^-\mat{H}^T\mat{S}^{-1}
    \end{equation}
    
    \item \textbf{Error state estimate}:
    \begin{equation}
        \delta\hat{\mat{x}} = \mat{K}\mat{y}
    \end{equation}
    
    \item \textbf{Apply corrections to nominal state}:
    \begin{align}
        \mat{r}^+ &= \mat{r}^- + \delta\mat{r} \\
        \mat{v}^+ &= \mat{v}^- + \delta\mat{v} \\
        \quat{q}^+ &= \delta\quat{q}(\vect{\theta}) \otimes \quat{q}^- \quad \text{(see Eq.~\ref{eq:error_quat}, \ref{eq:quat_correction})} \\
        \mat{b}_a^+ &= \mat{b}_a^- + \delta\mat{b}_a \\
        \mat{b}_g^+ &= \mat{b}_g^- + \delta\mat{b}_g
    \end{align}
    
    \item \textbf{Normalize quaternion}: $\quat{q}^+ \leftarrow \quat{q}^+ / \|\quat{q}^+\|$
    
    \item \textbf{Update covariance} (Joseph form for numerical stability):
    \begin{equation}
        \mat{P}^+ = (\mat{I} - \mat{K}\mat{H})\mat{P}^-(\mat{I} - \mat{K}\mat{H})^T + \mat{K}\mat{R}\mat{K}^T
    \end{equation}
    
    \item \textbf{Symmetrize}: $\mat{P}^+ \leftarrow \frac{1}{2}(\mat{P}^+ + \mat{P}^{+T})$
    
    \item \textbf{Reset error state} to zero (errors absorbed into nominal)
\end{enumerate}

%==============================================================================
\section{Implementation Checklist}
%==============================================================================

\begin{enumerate}
    \item[$\square$] Quaternion convention consistent (scalar-last, Hamilton)
    \item[$\square$] DCM formula matches quaternion convention (Eq.~\ref{eq:quat_to_dcm})
    \item[$\square$] Error quaternion composition order correct (Eq.~\ref{eq:quat_correction})
    \item[$\square$] Units converted: accel from $g$ to \si{\meter\per\second\squared}
    \item[$\square$] Units converted: gyro already in \si{\radian\per\second}
    \item[$\square$] Pressure converted to altitude for baro updates
    \item[$\square$] Quaternion normalized after each correction
    \item[$\square$] Covariance symmetrized after each update
    \item[$\square$] F matrix updated at each time step (depends on current state)
    \item[$\square$] G matrix updated at each time step (depends on $\mat{C}_b^n$)
\end{enumerate}

%==============================================================================
\section{State Dimension Summary}
%==============================================================================

\begin{center}
\begin{tabular}{lcc}
    \toprule
    & Nominal State & Error State \\
    \midrule
    Position & 3 & 3 \\
    Velocity & 3 & 3 \\
    Attitude & 4 (quaternion) & 3 (rotation vector) \\
    Accel bias & 3 & 3 \\
    Gyro bias & 3 & 3 \\
    \midrule
    \textbf{Total} & \textbf{16} & \textbf{15} \\
    \bottomrule
\end{tabular}
\end{center}

The dimension difference arises because:
\begin{itemize}
    \item The quaternion has 4 parameters but only 3 degrees of freedom (unit norm constraint)
    \item The error state uses a minimal 3-parameter rotation representation
    \item This is the standard error-state EKF approach for quaternion-based systems
\end{itemize}

%==============================================================================
\section{References}
%==============================================================================

\begin{enumerate}
    \item J.A. Farrell, ``Aided Navigation: GPS with High Rate Sensors,'' McGraw-Hill, 2008.
    \item J. Sol\`{a}, ``Quaternion kinematics for the error-state Kalman filter,'' arXiv:1711.02508, 2017.
    \item P.D. Groves, ``Principles of GNSS, Inertial, and Multisensor Integrated Navigation Systems,'' 2nd ed., Artech House, 2013.
    \item IEEE Std 952-1997, ``IEEE Standard Specification Format Guide and Test Procedure for Single-Axis Interferometric Fiber Optic Gyros.''
\end{enumerate}

\end{document}